%%%%%%%%%%  scan idx 269  |  folio 256  %%%%%%%%%%
\SGAPrintPage{256}%
that \(E^a/B^a\) is \(h\)-flat. By the criterion of flatness by fibers (EGA IV 11.3.10), we conclude that \(E^a/B^a\) is \(f\)-flat in a neighborhood of \(x\), q.e.d.

\sgaStableTarget{sga6.e03.cor.4-6-1}\textbf{Corollary 4.6.1.} \uline{Let} \(f:X\to S\) \uline{be a flat morphism locally of finite presentation}, \(x\in X\), \(s=f(x)\), \(E\in\Ob D^-(X)\), \([a,b]\) \uline{an interval of} \(\mathbb Z\). \uline{Denote by} \(i_s:X_s=X\times_S\Spec(k(s))\to X\) \uline{the canonical projection}, \uline{and put} \(E_{X_s}=\mathrm Li_s^*(E)\). \uline{Then we have}
\[
\begin{gathered}
 \parfamp_f(E)\subset[a,b]\text{ \uline{in a neighborhood of} }x\text{ \uline{and} }\tamp(E_{X_s})_x\subset[a,b]\\
 \Longleftrightarrow \parfamp(E)\subset[a,b]\text{ \uline{in a neighborhood of} }x.
\end{gathered}
\]
\uline{Consequently}, \(E\) \uline{is perfect if and only if} \(E\) \uline{is perfect rel.} \uline{to} \(f\) \uline{and} \(E_{X_s}\) \uline{is perfect for every} \(s\in S\).

\uline{Proof}. Apply (4.6) with \(X=Y\), \(f=\Id_X\).

\sgaStableTarget{sga6.e03.rem.4-6-2}\textbf{Remark 4.6.2.} From (4.6.1) and (I 5.10) it follows that, when \(f\) is locally of finite presentation, flat and has regular schemes as fibers, \(E\) perfect rel. to \(f\) is equivalent to \(E\) perfect (this equivalence had already been obtained (4.4) for \(f\) smooth).

\sgaStableTarget{sga6.e03.prop.4-7}\textbf{Proposition 4.7.} \uline{Let} \(S\) \uline{be a scheme}, \uline{and for} \(i=1,2\), \(f_i:X_i\to Y_i\) \uline{an} \(S\)-\uline{morphism locally of finite type}, \(E_i\in\Ob D(X_i)\), \([a_i,b_i]\) \uline{an interval of} \(\mathbb Z\). \uline{Suppose} \(X_1\) \uline{and} \(X_2\), \uline{as well as} \(Y_1\) \uline{and} \(Y_2\), \uline{are tor-independent over} \(S\) (1.5). \uline{Put} \(a=a_1+a_2\), \(b=b_1+b_2\), \(f=f_1\times_Sf_2\). \uline{Then we have}:
\begin{enumerate}[label=\alph*)]
\item \uline{If} \(\parfamp_{f_1}(E_1)\subset[a_1,b_1]\), \uline{and if} \(E_2\) \uline{is pseudo-coherent rel.} \uline{to} \(f_2\) \uline{and acyclic in degree outside} \([a_2,b_2]\), \(E_1\otimes_S^{\mathrm L}E_2\) (\uline{cf.} (1.6)) \uline{is pseudo-coherent rel.} \uline{to} \(f\) \uline{and acyclic in degree outside} \([a,b]\).
%%%%%%%%%%  scan idx 270  |  folio 257  %%%%%%%%%%
\SGAPrintPage{257}%
\item \(\parfamp_{f_i}(E_i)\subset[a_i,b_i]\) \uline{for} \(i=1,2\)
\[
 \Longrightarrow\parfamp_f(E_1\otimes_S^{\mathrm L}E_2)\subset[a,b].
\]
\end{enumerate}

\uline{Proof}. Combine (1.9) and (3.5).

\sgaStableTarget{sga6.e03.cor.4-7-1}\textbf{Corollary 4.7.1.} \uline{Under the hypotheses of} (4.7), \uline{if}, \uline{for} \(i=1,2\), \(f_i\) \uline{has perfect dimension} \(\leq n_i\), \(f\) \uline{has perfect dimension} \(\leq n_1+n_2\). \uline{Consequently}, \uline{if} \(f_1\) \uline{and} \(f_2\) \uline{are perfect}, \uline{the same holds for} \(f\).

\sgaStableTarget{sga6.e03.cor.4-7-2}\textbf{Corollary 4.7.2.} \uline{Let}
\[
\begin{tikzcd}
X \arrow[d,"f"'] & X' \arrow[l,"g"'] \arrow[d,"f'"]\\
Y & Y' \arrow[l,"h"']
\end{tikzcd}
\]
\uline{be a Cartesian square of schemes}, \uline{where} \(f\) \uline{is locally of finite type and} \(X\) \uline{and} \(Y'\) \uline{are supposed tor-independent over} \(Y\), \uline{and let} \(E\in\Ob D(X)\), \uline{and} \([a,b]\) \uline{be an interval of} \(\mathbb Z\). \uline{We have}:
\[
 \parfamp_f(E)\subset[a,b]\Longrightarrow\parfamp_{f'}(\mathrm Lg^*(E))\subset[a,b],
\]
\uline{and consequently the functor} \(\mathrm Lg^*:D^-(X)\to D^-(X')\) \uline{sends} \(D(f)_{\parf}\) (resp. \(D(f)_{\uparf}\)) into \(D(f')_{\parf}\) (resp. \(D(f')_{\uparf}\)) (4.2).

\sgaStableTarget{sga6.e03.prop.4-8}\textbf{Proposition 4.8.} \uline{Let} \(S\) \uline{be a scheme}, \(f:X\to Y\) \uline{a proper} \(S\)-\uline{morphism of} \(S\)-\uline{schemes locally of finite type}, \(E\in\Ob D^b(X)\). \uline{Then}, \uline{subject to the validity of} (2.1) \uline{for} \(f\) (\uline{hence in any case when} \(S\) \uline{is locally noetherian} \uline{or when} \(f\) \uline{is projective}):
\begin{enumerate}[label=\alph*)]
\item \uline{If} \(E\) \uline{is perfect rel.} \uline{to} \(S\), \uline{the same holds for} \(Rf_*(E)\).
\item \uline{Suppose} \(S\) \uline{quasi-compact}, \(X\) \uline{and} \(Y\) \uline{of finite type over} \(S\); \uline{let} \([a,b]\) \uline{be an interval of} \(\mathbb Z\), \uline{and} \(N\) \uline{an integer such that} \(R^if_*(M)=0\) \uline{for every} \(i>N\) \uline{and}
%%%%%%%%%%  scan idx 271  |  folio 258  %%%%%%%%%%
\SGAPrintPage{258}%
\uline{every} \(\underline{O}_X\)-\uline{Module} \(M\) \uline{that is quasi-coherent} (\uline{such an integer exists by} (EGA III 1.4.12 \uline{and} IV 1.7.21))\,; \uline{we then have}
\[
 \parfamp_S(E)\subset[a,b]\Longrightarrow\parfamp_S(Rf_*(E))\subset[a,b+N].
\]
\end{enumerate}
\uline{Consequently}, \uline{the functor} \(Rf_*\) \uline{sends} \(D(X)_{S\text{-}\uparf}\) \uline{into} \(D(Y)_{S\text{-}\uparf}\)\SGAAuthorFootnote{1}{moreover, \(D(X)_{S\text{-}\uparf}=D(X)_{S\text{-}\parf}\) (and likewise with \(Y\)), since \(X\) and \(Y\) are quasi-compact (cf. (4.2)).}.

\uline{Proof}. Assertion a) follows from b), and b) is obtained by combining (2.1) and (3.7.1).

\sgaStableTarget{sga6.e03.cor.4-8-1}\textbf{Corollary 4.8.1.} \uline{Let} \(f:X\to Y\) \uline{be a proper and perfect morphism}, \uline{and} \(E\in\Ob D^+(X)\). \uline{Then}, \uline{subject to the validity of} (2.1) \uline{for} \(f\), \uline{if} \(E\) \uline{is perfect}, \uline{the same holds for} \(Rf_*(E)\)\,; \uline{moreover}, \uline{if} \(Y\) \uline{is quasi-compact}, \uline{the functor} \(Rf_*\) \uline{sends} \(D(X)_{\uparf}\) \uline{into} \(D(Y)_{\uparf}\)\SGAAuthorFootnote{2}{We have \(D(X)_{\uparf}=D(X)_{\parf}\) (and likewise with \(Y\)).}.

\uline{Proof}. Indeed, when \(E\) is perfect, \(E\) is perfect rel. to \(f\) by (4.5.1 a)).

\sgaStableTarget{sga6.e03.prop.4-9}\textbf{Proposition 4.9}\SGAAuthorFootnote{3}{This result and its corollaries will not be used in the remainder of this Seminar, but play an important role in the Lefschetz--Verdier formula for coherent sheaves; see SGA 5 Exp.\ III.}. \uline{Let} \(S\) \uline{be a noetherian scheme}, \(X\) \uline{and} \(Y\) \uline{be} \(S\)-\uline{schemes of finite type}, \(f:X\to Y\) \uline{a compactifiable} \(S\)-\uline{morphism} (\uline{i.e.} \uline{a composite of an open immersion and a proper morphism}), \(E\in\Ob D(X)\), \(F\in\Ob D^+(Y)\). \uline{If} \(E\) \uline{is perfect rel.} \uline{to} \(f\), \uline{and} \(F\) \uline{perfect rel.} \uline{to} \(S\), \(\mathrm R\underline{\Hom}(E,f^!(F))\) \uline{is perfect rel.} \uline{to} \(S\)\SGAAuthorFootnote{4}{For the definition of the functor \(f^!\), see [H] App.\ n\textsuperscript{o}~4.}.

\uline{Proof}. The question being local on \(X\), we may suppose that \(f\) factors as \(X\xrightarrow{i}X'\xrightarrow{f'}Y\), where \(i\) is a closed immersion and \(f'\) a smooth morphism. When \(E\) is perfect rel. to \(f\), \(E\) is perfect rel. to \(i\) (4.4), and when \(F\) is perfect rel. to \(S\), \(f'^!(F)\) is perfect rel. to \(S\) by (4.5.1 b)), since \(f'^!(F)\) is locally isomorphic to \(f'^*(F)\) up to a shift. Since \(f^!\simeq i^!f'^!\), we therefore reduce to proving the assertion when \(f\) is a closed immersion. Thus suppose \(E\) perfect rel. to \(f\), i.e. \(f_*(E)\) perfect, and \(F\) perfect rel. to \(S\). We must show that \(f_*\mathrm R\underline{\Hom}(E,f^!F)\) is
%%%%%%%%%%  scan idx 272  |  folio 259  %%%%%%%%%%
\SGAPrintPage{259}%
perfect rel. to \(S\). But by the duality theorem ([H] III 6.7), we have
\[
 f_*\mathrm R\underline{\Hom}(E,f^!F)\simeq\mathrm R\underline{\Hom}(f_*(E),F).
\]
Since \(f_*(E)\) is perfect, we have, by (I 7.7),
\[
 \mathrm R\underline{\Hom}(f_*(E),F)\simeq(f_*(E))^{\vee}\otimes^{\mathrm L}F,
\]
and (4.5) (where \(X=Y\) and \(f=\Id_X\)) gives the result.

\sgaStableTarget{sga6.e03.cor.4-9-1}\textbf{Corollary 4.9.1.} \uline{Under the hypotheses of} (4.9), \uline{if} \(f\) \uline{is perfect}, \(f^!\) \uline{sends} \(D(Y)_{S\text{-}\parf}\) \uline{into} \(D(X)_{S\text{-}\parf}\).

\uline{Proof}. Take \(E=\underline{O}_X\).

\sgaStableTarget{sga6.e03.cor.4-9-2}\textbf{Corollary 4.9.2.} \uline{Let} \(S\) \uline{be a noetherian scheme}, \(f:X\to S\) \uline{a compactifiable morphism}. \uline{Put} \(f^!(\underline{O}_S)=K_X\), \(\mathrm R\underline{\Hom}(.,K_X)=\underline{D}_X\). \uline{For} \(E\in\Ob D(f)_{\parf}\), \uline{we have} \(\underline{D}_X(E)\in\Ob D(f)_{\parf}\), \uline{and the canonical arrow} (I 7.2)
\[
 E\longrightarrow \underline{D}_X\underline{D}_X(E)
\]
\uline{is an isomorphism}.

\uline{Proof}. For the first assertion, apply (4.9) with \(Y=S\), \(F=\underline{O}_S\). For the second, we easily reduce to the case where \(f\) is a closed immersion. It then suffices to prove that the arrow obtained from the canonical arrow by applying \(f_*\) is an isomorphism. By global duality ([H] III 6.7), we have
\[
 f_*\underline{D}_X\simeq\underline{D}_Sf_*.
\]
We therefore reduce to showing that the canonical arrow
\[
 f_*(E)\longrightarrow \underline{D}_S\underline{D}_Sf_*(E)
\]
is an isomorphism for \(E\) perfect rel. to \(f\), i.e. \(f_*(E)\) perfect, which follows from ordinary biduality for perfect complexes (I 7.2).

\sgaStableTarget{sga6.e03.sec.5}\section*{5. Applications: exchange and semi-continuity theorems\SGAAuthorFootnote{1}{The results of this section will not be used in the remainder of the Seminar.}}

\sgaStableTarget{sga6.e03.lem.5-1}\textbf{Lemma 5.1.} \uline{Let} \((X,\underline{O}_X)\) \uline{be a ringed topos} (\uline{commutative}), \uline{denote by} \(\Mod(X)\) \uline{the category of} \(\underline{O}_X\)-\uline{Modules}, \(D(X)=D(\Mod(X))\). \uline{Let} \(E\in\Ob D^-(X)\), \uline{and} \(p\in\mathbb Z\). \uline{Consider the following conditions}:
%%%%%%%%%%  scan idx 273  |  folio 260  %%%%%%%%%%
\SGAPrintPage{260}%
\begin{enumerate}[label=\alph*)]
\item \uline{The functor} \(M\mapsto H^p(E\otimes^{\mathrm L}M)\) \uline{from} \(\Mod(X)\) \uline{to} \(\Mod(X)\) \uline{is right exact}.
\item[a')] \uline{The canonical arrow}
\[
 H^p(E)\otimes M\longrightarrow H^p(E\otimes^{\mathrm L}M)
\]
\uline{is an isomorphism for every} \(M\in\Ob\Mod(X)\).
\item[b)] \uline{The functor} \(M\mapsto H^{p+1}(E\otimes^{\mathrm L}M)\) \uline{from} \(\Mod(X)\) \uline{to} \(\Mod(X)\) \uline{is left exact}.
\item[b')] \uline{The canonical arrow defined below}
\begin{equation}
 H^{p+1}(E\otimes^{\mathrm L}M)\longrightarrow\underline{\Hom}(H^{-p-1}(E^{\vee}),M)\tag{$*$}
\end{equation}
\uline{is an isomorphism for every} \(M\in\Ob\Mod(X)\).
\end{enumerate}
\uline{Then we have}:
\[
 \mathrm a)\Longleftrightarrow\mathrm a')\Longleftrightarrow\mathrm b)\Longleftarrow\mathrm b').
\]
\uline{Moreover}, \uline{if} \(E\) \uline{is pseudo-coherent}, \uline{the conditions} a), a'), b), b') \uline{are equivalent}.

\uline{Definition of} \((*)\): The canonical arrow (I 7.5) \(E^{\vee}\otimes^{\mathrm L}E\longrightarrow\underline{O}_X\), where \(E^{\vee}=\mathrm R\underline{\Hom}(E,\underline{O}_X)\), gives, by tensoring with \(M\in\Ob\Mod(X)\), an arrow \(E^{\vee}\otimes^{\mathrm L}E\otimes^{\mathrm L}M\longrightarrow M\); hence, by the adjunction formula dear to Cartan (I 7.4), a canonical arrow
\begin{equation}
 E\otimes^{\mathrm L}M\longrightarrow\mathrm R\underline{\Hom}(E^{\vee},M).\tag{$**$}
\end{equation}
Composing the arrow induced by \((**)\) on \(H^{p+1}\) with the canonical arrow
\begin{equation}
 \underline{\Ext}^{p+1}(E^{\vee},M)\longrightarrow\underline{\Hom}(H^{-p-1}(E^{\vee}),M),\tag{$***$}
\end{equation}
we obtain the arrow \((*)\).

\uline{Proof}. The implications
\[
 \mathrm a')\Longrightarrow\mathrm a)\Longleftrightarrow\mathrm b)\Longleftarrow\mathrm b')
\]
are trivial. Let us prove \(\mathrm a)\Longrightarrow\mathrm a')\). Let \(M\) be an \(\underline{O}_X\)-Module. There is an exact sequence \(M^1\longrightarrow M^0\longrightarrow M\longrightarrow0\), where \(M^0\) and \(M^1\) are direct sums of sheaves
